\chapter{QCD and the Parton Idea: Theoretical Foundations}

Everything in this book rests on three discoveries of the early 1970s:
that the strong interaction becomes weak at short distances, that parton
distributions obey renormalization-group evolution, and that the theory
can be regularized on a spacetime lattice without losing gauge invariance.
The library cites all three whenever it needs to say why lattice QCD is
worth the effort.

% ----------------------------------------------------------------------
\section{Asymptotic Freedom}
\papercard{D.\ J.\ Gross and F.\ Wilczek}{Ultraviolet behavior of non-abelian gauge theories}
{Phys.\ Rev.\ Lett.\ \textbf{30}, 1343 (1973); companion paper: H.\ D.\ Politzer,
\emph{Reliable perturbative results for strong interactions?},
Phys.\ Rev.\ Lett.\ \textbf{30}, 1346 (1973)}
{No arXiv (preprint era) | DOI: 10.1103/PhysRevLett.30.1343 / .1346}
\whyselected{The discovery that non-abelian gauge theory is asymptotically
free is the single fact that makes QCD a calculable theory and parton
distributions meaningful objects. It is cited in the library wherever the
reason for doing lattice QCD is stated (e.g.\ the BMW spectrum and
neutron--proton mass-difference papers, and the lattice-QCD thesis in the
collection).}
\physics{In Yang--Mills theory with $n_f$ quark flavors the one-loop
beta function is negative,
\begin{equation}
  \beta(g) = -\frac{g^3}{16\pi^2}\Big(11 - \frac{2}{3}n_f\Big) + \mathcal{O}(g^5),
\end{equation}
so the coupling decreases logarithmically at short distance,
$\alpha_s(Q^2)\to 0$ as $Q^2\to\infty$. Deep-inelastic scattering then
sees nearly free quarks (Bjorken scaling to leading order), while at large
distances the coupling grows --- confinement is not proven here but made
plausible as the same phenomenon viewed from the infrared side. The sign
of $\beta_0$ is fixed by the gluon self-interaction: only non-abelian
gauge theories screen antiscreened.}
\impact{Every global PDF analysis in the library (CT18, NNPDF, CJ, JAM)
evolves its partons with splitting functions whose existence is guaranteed
by this result; every lattice paper invokes it when motivating QCD as the
theory to simulate. It is the root of the entire tree.}
\cardend

% ----------------------------------------------------------------------
\section{DGLAP Evolution}
\papercard{G.\ Altarelli and G.\ Parisi}
{Asymptotic freedom in parton language}
{Nucl.\ Phys.\ B \textbf{126}, 298 (1977)}
{No arXiv (preprint era)}
\whyselected{The DGLAP equation translated the operator-product-expansion
view of scaling violations into the parton language in which every PDF
paper of this library is written. It appears in the reference lists of the
Ji~2013 founding paper, the pseudo-PDF papers of Radyushkin and Orginos
et al., the gradient-flow PDF papers, and many more ($\sim$8/99).}
\physics{Quark distributions are not constants but scale-dependent
amplitudes obeying an evolution equation of probabilistic form,
\begin{equation}
  \frac{\partial q(x,\mu)}{\partial\ln\mu^2}
    = \frac{\alpha_s(\mu)}{2\pi}\,(P_{qq}\otimes q + P_{qg}\otimes g),
  \qquad
  (P\otimes q)(x)=\int_x^1\frac{dz}{z}\,P(z)\,q\!\left(\frac{x}{z},\mu\right).
\end{equation}
The splitting functions $P_{ij}$ are computed from collinear factorization;
in Mellin moment space they become matrices acting on moments
$M_n=\int dx\,x^{n-1}q(x)$, which is exactly the language of lattice
moments calculations. DGLAP is the ``vertical'' axis of every plot in this
book's subject area: any distribution extracted on the lattice at some
scale $\mu$ must be evolved before comparison with phenomenology.}
\impact{The matching kernels of quasi- and pseudo-PDFs reproduce the
one-loop $P_{ij}$ in their small-$z$ limit --- this consistency check
appears verbatim in Ji~2013, Radyushkin~2017, Balitsky--Morris--Radyushkin
2020 and the smeared-quasi-PDF papers. Modern three-loop splitting
functions (Moch--Vermaseren--Vogt 2004) underpin the NNLO fits used as
benchmarks throughout the library.}
\cardend

% ----------------------------------------------------------------------
\section{Lattice Gauge Theory}
\papercard{Kenneth G.\ Wilson}
{Confinement of quarks}
{Phys.\ Rev.\ D \textbf{10}, 2445--2459 (1974)}
{No arXiv (preprint era) | DOI: 10.1103/PhysRevD.10.2445}
\whyselected{The founding document of lattice field theory and the
conceptual ancestor of every calculation in the library. Within this
corpus it is cited by the BMW mass papers, the AI-sampling papers, the
agentic-computing papers, and the chiral-phase-transition thesis
($\sim$4/99 directly; effectively universal). A full English LaTeX
transcription of this paper exists in the companion \texttt{books/}
collection together with its Chinese translation.}
\physics{Wilson replaced continuum spacetime by a hypercubic lattice while
keeping gauge invariance exact: the degrees of freedom are link variables
$U_\mu(n)\in SU(3)$, parallel transports along links, and the action is
built from plaquettes,
\begin{equation}
  S_W = \frac{\beta}{3}\sum_{n,\mu<\nu}\Big(1-\frac{1}{3}\mathrm{Re\,Tr}\,U_{\mu\nu}(n)\Big),
  \qquad \beta = \frac{2N_c}{g^2}.
\end{equation}
Gauge invariant observables are traces of products of links along closed
paths --- Wilson loops. The area law $\langle W(C)\rangle\sim e^{-\sigma A}$
signals confinement through a nonzero string tension $\sigma$, and the
static quark potential rises linearly. Strong-coupling expansion proves
the area law at $\beta=0$; asymptotic freedom guarantees the Coulomb law
at $\beta\to\infty$: both regimes live on one lattice, connected by the
crossover that Monte Carlo (Chapter~2) was invented to explore.}
\impact{Every ensemble, every correlator, every flow time in the library's
99 papers exists because Wilson wrote this paper. His name also marks the
Wilson line --- the path-ordered exponential of links --- which will turn
out to be precisely the object whose renormalization (Chapters~5--6)
decides whether parton physics can be extracted from the lattice at all.}
\cardend
